\section{Discussion and Limitations}

\paragraph{Interpretation of the stopping rule.}
The theorem predicts when optimization reaches the unique-data frontier, not
the exact validation optimum in every noisy application. Label noise can create
an earlier variance- or overfitting-limited optimum, so the coverage horizon is
best viewed as an upper limit on useful unregularized reuse.

\paragraph{Known versus estimated frequencies.}
The sharp mask uses population $p_j$ values; a plug-in mask needs uniform
small-count control. The plain output needs neither a mask nor an abort event,
at the cost of the logarithm in our present proof.

\paragraph{Representation and algorithmic scope.}
The theorem assumes independent addressable sparse coordinates, clipped
with-replacement replay, averaging, noiseless labels, and
$a-s+1<b<a+1$. Fixed or learned features, cyclic replay, reshuffling, and last
iterates require additional analysis. Progressive subsampling is an
achievability construction, not an optimality claim.

\paragraph{Practical implication.}
Input occupancy and conditional second moments indicate when to redirect
resources from replay to unique data or model capacity.

\section{Conclusion}

Sparse data impose a resource missing from total-update counts: unique feature
coverage. We derive its unavoidable floor, characterize its interaction with
capacity and optimization, and prove stable coactive replay reaches it at the
predicted epoch scale. Width sweeps and an input-only estimator validate all
three frontiers. Replay while optimization dominates; after coverage takes over,
spend on unique data or model capacity.

